\section{Experiments}
\label{sec:experiments}

\begin{figure*}[t]
  \centering
  \includegraphics[width=\linewidth]{neurolink_v5_1_dashboard.png}
  \caption{%
    \textbf{NeuroLink 3.0 v5.1 system dashboard over 435 deployment days.}
    \textbf{Top row:} 17-month Longitudinal Drift Index (LDI) timeline with
    7-day rolling mean (blue), bootstrap 95\% CI band (orange shading),
    HIGH anomaly days (red circles, LDI $\geq\theta_{\text{hi}}$), and
    MARGINAL days (orange squares, $\theta_{\text{lo}}\leq$LDI$<\theta_{\text{hi}}$).
    \textbf{Middle left:} Training and validation loss convergence; early stopping
    fires at epoch 54.
    \textbf{Middle right:} Per-feature validation $R^2$ — 10 of 12 features
    exceed 0.90; all exceed 0.85.
    \textbf{Bottom row (left to right):} Circadian drift gate showing mean LDI
    by UTC hour, monthly LDI heatmap, hand asymmetry index over time,
    and a 2-D PCA projection of the 10-dimensional Muscle Memory Signature
    latent space coloured by WPM.%
  }
  \label{fig:dashboard}
\end{figure*}

\subsection{Experimental Setup}
\label{sec:setup}

The 42,501 lessons are split chronologically: the earliest 85\% (lessons
1--36,125) form the training set, the next 7.5\% the validation set, and the
final 7.5\% the held-out test set.
This split preserves temporal order, preventing any future information from
leaking into training.
The first 36,125 training lessons also define the ``normal'' baseline from
which the bootstrap anomaly thresholds are derived.

All features are standardized with zero mean and unit variance using
statistics computed \emph{only} on the training split, then applied
identically to validation and test splits.
The model is implemented in PyTorch 2.x and trained on an NVIDIA T4 GPU
(Google Colab) with batch size 64 for up to 200 epochs.

\subsection{Ablation Study}
\label{sec:ablation}

\Cref{tab:ablation} traces the progression from the baseline v4 system to
the final v5.1 model, isolating the contribution of each design choice.

\begin{table}[h]
  \centering
  \caption{Ablation study: cumulative improvements from v4 to v5.1.
           All metrics evaluated on the held-out test split.
           \textbf{Bold} marks the best value per column.}
  \label{tab:ablation}
  \resizebox{\columnwidth}{!}{%
  \begin{tabular}{lcccc}
    \toprule
    \textbf{Variant} & \textbf{Val $R^2$} & \textbf{EV $R^2$} & \textbf{ESRP $R^2$} & \textbf{Best Ep.} \\
    \midrule
    v4 (rolling-mean EV/ESRP)    & 0.9502 & 0.234 & 0.113 & 47 \\
    v5 (residual + LR 5e-3)      & 0.8860 & 0.948 & 0.730 & 10 \\
    v5.1 (+ ReduceLR + weights)  & \textbf{0.9586} & \textbf{0.954} & \textbf{0.854} & \textbf{54} \\
    \bottomrule
  \end{tabular}}
\end{table}

\paragraph{v4 → v5: Residual Reformulation.}
Replacing rolling-mean EV and ESRP with their residual formulations
(Eqs.~\ref{eq:ev}--\ref{eq:esrp}) raises EV $R^2$ from 0.234 to 0.948
and ESRP $R^2$ from 0.113 to 0.730.
However, v5 converges prematurely at epoch 10 due to the initial LR of
$5\times10^{-3}$: residual features are structurally easy to predict,
so the loss drops quickly while harder features (\texttt{lesson\_dur\_z},
\texttt{practice\_gap}) are insufficiently optimized.

\paragraph{v5 → v5.1: ReduceLROnPlateau + Weight Rebalancing.}
Lowering the initial LR to $2\times10^{-3}$ and switching to
\texttt{ReduceLROnPlateau} (factor 0.5, patience 5) allows the optimizer
to escape early plateaus.
Rebalancing feature weights — raising \texttt{lesson\_dur\_z} from 2.0 to
4.0 and \texttt{practice\_gap} from 1.5 to 2.5, while lowering EV and
ESRP from 2.0 to 1.0 — redirects model capacity toward harder features
without sacrificing easy ones.
The model trains to epoch 54 and uses nearly the full DP budget
($\varepsilon{=}7.9983$), yielding the best overall $R^2$ across all variants.

\subsection{Per-Feature Reconstruction Quality}
\label{sec:perfeat}

\Cref{tab:perfeat} reports the per-feature validation $R^2$ for the final
v5.1 model.
All 12 features exceed $R^2{=}0.82$, and 10 of 12 exceed 0.90.

\begin{table}[h]
  \centering
  \caption{Per-feature validation $R^2$ for NeuroLink 3.0 v5.1.}
  \label{tab:perfeat}
  \resizebox{\columnwidth}{!}{%
  \begin{tabular}{lcc}
    \toprule
    \textbf{Feature} & \textbf{$R^2$} & \textbf{Weight $w_j$} \\
    \midrule
    EV (residual)          & 0.954 & 1.0 \\
    RSV                    & 0.964 & 2.0 \\
    SSF                    & 0.980 & 1.0 \\
    ESRP (residual)        & 0.854 & 1.0 \\
    CTA (sin)              & 0.940 & 1.0 \\
    wpm\_vs\_50roll        & 0.890 & 1.0 \\
    CTA\_cos               & 0.981 & 1.0 \\
    lesson\_dur\_z         & 0.827 & 4.0 \\
    active\_char\_frac     & 0.954 & 1.0 \\
    hand\_asymmetry        & 0.990 & 2.0 \\
    practice\_gap          & 0.904 & 2.5 \\
    error\_entropy         & 0.989 & 2.0 \\
    \midrule
    \textbf{Variance-weighted mean} & \textbf{0.9586} & --- \\
    \bottomrule
  \end{tabular}}
\end{table}

The two lowest-scoring features (\texttt{lesson\_dur\_z}, $R^2{=}0.827$, and
\texttt{ESRP}, $R^2{=}0.854$) both reflect high-variance, heavy-tailed
distributions that are intrinsically harder to compress.
The hand asymmetry and error entropy features, despite their 2.0 weights,
achieve near-perfect reconstruction ($\geq0.989$), suggesting they encode
highly regular within-user patterns.

\subsection{Longitudinal Drift Index Analysis}
\label{sec:ldi_results}

\Cref{tab:ldi} summarizes LDI statistics and anomaly tier counts over the
full 435-day deployment period.
\Cref{fig:dashboard} (top row) shows the complete timeline.

\begin{table}[h]
  \centering
  \caption{LDI summary and anomaly tier distribution over 435 deployment days.}
  \label{tab:ldi}
  \begin{tabular}{lc}
    \toprule
    \textbf{Metric} & \textbf{Value} \\
    \midrule
    Training LDI mean              & 0.0548 \\
    Training LDI std               & 0.0758 \\
    Bootstrap $\theta_{\text{lo}}$ & 0.1744 \\
    Bootstrap $\theta_{\text{hi}}$ & 0.2352 \\
    \midrule
    NORMAL days  ($< \theta_{\text{lo}}$)                              & 399 (91.7\%) \\
    MARGINAL days ($\theta_{\text{lo}} \leq \cdot < \theta_{\text{hi}}$) & 21 (4.8\%) \\
    HIGH days ($\geq \theta_{\text{hi}}$)                              & 15 (3.4\%) \\
    \bottomrule
  \end{tabular}
\end{table}

The system achieves a low false-positive rate: 91.7\% of days are classified
as NORMAL, consistent with the assumption that most days reflect the
user's established behavioral baseline.
HIGH anomalies (3.4\%) cluster around days with irregular practice schedules
(high \texttt{practice\_gap}) or atypical session durations, which the model
correctly identifies as deviating from baseline.
The MARGINAL tier (4.8\%) provides a warning band that captures days
near-threshold without triggering the HIGH alert, reducing unnecessary alarms.

\subsection{Ghost V3 Drift-Adjusted Predictor}
\label{sec:ghost}

The \emph{Empirical Ghost V3} predictor provides a closed-form estimate of
expected WPM as a function of observed LDI deviation from baseline.
We fit a linear regression across all 42,501 lessons:
\[
  \widehat{\text{WPM}} = a + b \cdot \text{ldi\_ratio}
\]
where $\text{ldi\_ratio} = \text{LDI}_t / \overline{\text{LDI}}_{\text{train}}$
is the ratio of current daily LDI to the training-period mean.
The fitted coefficients yield an empirical scale factor
$\gamma = -b / \overline{\text{WPM}}$, which replaces the hand-tuned constant
used in Ghost V2.

The key distinction from prior Ghost versions is that $\gamma$ is derived
entirely from data: it reflects how much WPM typically drops for each unit
of normalized LDI increase, measured over the full longitudinal record.
This eliminates an unverified hyperparameter and makes the predictor
falsifiable: $\gamma$ can be re-estimated as more data accumulates without
retraining the autoencoder.

\subsection{Privacy Accounting}
\label{sec:privacy}

Training with Opacus DP-SGD (RDP accountant, $\delta{=}10^{-5}$, $C{=}1.0$,
$\sigma{=}0.5994$, batch size 64) over 54 epochs yields a final privacy
expenditure of $\varepsilon{=}7.9983$, within 0.0017 of the target $\varepsilon{=}8.0$.
This guarantees that the trained model cannot output significantly different
results on two datasets differing by a single lesson, bounding the
information any single training lesson can contribute to the behavioral baseline.
